%% ma: L11-B14-0003
%% lop: 11
%% chuong: 4
%% bai: 14
%% dang: 11.14.1
%% loai: TL
%% muc: VD
%% dap_an: "-"
%% nguon: "tai-lieu-hoc-tap-mon-toan-lop-11-theo-chuong-trinh-giao-duc-pho-thong-2018.pdf (Lớp Toán Thầy Duy, Phạm Lê Duy, THPT Chu Văn An), Chương 4 Bài 5, Dạng 3, Ví dụ 3 (tr.523)"
%% kien_thuc: "Lớp 11 Bài 14 (phép chiếu song song)"
%% trang_thai: cho-duyet
%% ly_do: "Dạng mới đề xuất (11.14.x), chờ giáo viên duyệt dạng; tài liệu gốc không có lời giải mẫu, đáp án do kho tự giải và kiểm bằng sympy."
%% ghi_chu: "Đề gốc không có lời giải mẫu; kho giải ra hình chiếu suy biến thành đoạn thẳng MN (mặt phẳng (C'MN) song song với phương chiếu AO'). Đã kiểm bằng tọa độ."
\begin{ex}
Cho hình hộp chữ nhật $ABCD.A^{\prime}B^{\prime}C^{\prime}D^{\prime}$. Gọi $O=AC\cap BD$ và $O^{\prime}=A^{\prime}C^{\prime}\cap B^{\prime}D^{\prime}$. Điểm $M,N$ lần lượt là trung điểm của $AB$ và $CD$. Qua phép chiếu song song theo phương $AO^{\prime}$ lên mặt phẳng $(ABCD)$, xác định hình chiếu của tam giác $C^{\prime}MN$.
\loigiai{\sloppy $M,N\in(ABCD)$ nên hình chiếu của chúng là chính $M,N$. Tứ giác $AOC^{\prime}O^{\prime}$ có $AO\parallel O^{\prime}C^{\prime}$ và $AO=O^{\prime}C^{\prime}$ (cùng bằng nửa đường chéo của hai đáy bằng nhau) nên là hình bình hành, suy ra $OC^{\prime}\parallel AO^{\prime}$; vậy hình chiếu của $C^{\prime}$ theo phương $AO^{\prime}$ lên $(ABCD)$ là $O$. Vì $O$ là tâm hình chữ nhật $ABCD$ và $M,N$ là trung điểm của $AB,CD$ nên $O$ là trung điểm của $MN$, tức $O,M,N$ thẳng hàng. Vậy hình chiếu của $\triangle C^{\prime}MN$ là đoạn thẳng $MN$.}
\end{ex}
